% Irodov Part V §5.4, original problems 158--175.
% Source: supplied PDF pp.222--224 (printed pp.213--215); answers pp.341--342.
% Draft for scholarly checking. All polarization intensity formulae refer
% to energy fluxes; clarify normal-to-ray vs normal-to-interface intensity.

\polarizationproblem{158}{213}
强度为 $I_0$ 的自然单色平面波垂直入射不透明屏；
圆孔恰对应观察点 $P$ 的第一菲涅耳半波带。
用两块透振轴正交、相接无缝的理想偏振片盖住圆孔。
求分界线 (a) 为圆孔的一条直径，
(b) 为将第一菲涅耳半波带按光程差等分为两半的同心圆周时，$P$ 点的光强。
\exsolution
令完全敞开波前在 $P$ 处的复振幅为 $A_\infty$。
菲涅耳矢量图表明完整第一半波带振幅为 $2A_\infty$。
自然光的两个正交统计分量各有 $I_0/2$；
两块偏振片输出始终正交，故只相加强度。
(a) 对称的两个半圆各给全波带振幅的一半，即 $A_\infty$，
所以
\[
 I_P=2\cdot\frac{I_0}{2}
 \left|\frac{A_\infty}{A_\infty}\right|^2
 =\boxed{I_0}.
\]
(b) 以菲涅耳相位 $\psi$ 积分，第一半波带对应
$0\le\psi\le\pi$；两部分分别是 $[0,\pi/2]$ 与
$[\pi/2,\pi]$。各自振幅与完整第一半波带振幅之比为
\[
 \frac{|e^{i\pi/2}-1|}{|e^{i\pi}-1|}=\frac1{\sqrt2},
\]
因此两部分各自相对全开放场的振幅模为 $\sqrt2$，
\[
 I_P=2\cdot\frac{I_0}{2}\cdot 2
 =\boxed{2I_0}.
\]
\excheck 分界圆周是\emph{半个菲涅耳波带}的相位界，
不能把半个面积与半个复振幅混同；(a)(b) 与书末结果一致。

\polarizationproblem{159}{214}
入射线偏振光的能流率为
$\Phi_0=4.0\,\mathrm{mW}$，理想检偏器绕光束轴以
$\omega=21\,\mathrm{rad/s}$ 匀速旋转。
求检偏器转一周内透射的光能。
\exsolution
在时间 $t$，透振轴与入射偏振轴夹角为
$\omega t+\phi_0$。马吕斯定律给
$\Phi(t)=\Phi_0\cos^2(\omega t+\phi_0)$。
一周历时 $T=2\pi/\omega$，
\[
 E=\int_0^T\Phi(t)\,dt
 =\frac{\Phi_0 T}{2}
 =\boxed{\frac{\pi\Phi_0}{\omega}
 \simeq 5.98\times10^{-4}\,\mathrm{J}=0.60\,\mathrm{mJ}}.
\]
\excheck 周期平均透射率为 $1/2$，结果与初相位无关。

\polarizationproblem{160}{214}
自然光依次经过 $N=6$ 块理想尼科耳棱镜，
每后一块的透振方向比前一块转过
$\phi=30^\circ$。求总能流透射比。
\exsolution
第一块将自然光强度减半；其后每跨越相邻棱镜按
马吕斯定律乘上 $\cos^2\phi$。所以
\[
 \boxed{\eta=\frac12\cos^{2(N-1)}\phi
 =\frac12\left(\frac34\right)^5
 =\frac{243}{2048}\simeq0.119}.
\]
\excheck 指数是 $2(N-1)=10$ 而非 $2N$。

\polarizationproblem{161}{214}
自然光经过三块相同且共轴的实际偏振片；
中间偏振片的透振方向分别与两侧偏振片成
$\phi=60^\circ$，即两侧的透振方向互相平行。
每片对沿透振轴的线偏振光的最大透射率为
$\tau=0.81$。问光强降低为原来的多少分之一。
\exsolution
每片的强度损耗因子为 $\tau$；
第一片后的光强为 $I_0\tau/2$。
通过第二片与第三片时各乘 $\tau\cos^2\phi$，故
\[
 I=\frac{I_0}{2}\tau^3\cos^4\phi,\qquad
 \boxed{\frac{I_0}{I}=\frac{2}{\tau^3\cos^4\phi}
 \simeq60.2}.
\]
\excheck $\tau$ 是强度透射率而非振幅透射率，不得再次平方。

\polarizationproblem{162}{214}
部分偏振光的偏振度 $P=0.25$，
求其中完全偏振成分与自然光成分的强度比。
\exsolution
记 $I=I_{\mathrm p}+I_{\mathrm n}$。
在检偏器旋转中
$I_{\max}=I_{\mathrm p}+I_{\mathrm n}/2$，
$I_{\min}=I_{\mathrm n}/2$。
因此
\[
 P=\frac{I_{\max}-I_{\min}}{I_{\max}+I_{\min}}
 =\frac{I_{\mathrm p}}{I_{\mathrm p}+I_{\mathrm n}},
 \qquad
 \boxed{\frac{I_{\mathrm p}}{I_{\mathrm n}}=
 \frac{P}{1-P}=\frac13}.
\]
\excheck 偏振度的分母是入射\emph{总}强度。

\polarizationproblem{163}{214}
部分线偏振光经尼科耳棱镜检偏；
棱镜自最大透射位置旋转 $\phi=60^\circ$，
透射强度降为原来的 $1/\eta$，其中 $\eta=3.0$。
求入射光偏振度。
\exsolution
将入射光分为自然强度 $I_{\mathrm n}$ 和线偏振强度
$I_{\mathrm p}$，设其偏振度为 $P$。检偏器的透射率曲线可写成
\[
 I(\theta)=\frac I2(1+P\cos2\theta).
\]
因此
\[
 \eta=\frac{I(0)}{I(\phi)}
 =\frac{1+P}{1+P\cos2\phi},\quad
 \boxed{P=\frac{\eta-1}{1-\eta\cos2\phi}=0.80}.
\]
\excheck $0\le P\le1$；零偏振时旋转检偏器不应改变光强。

\polarizationproblem{164}{214}
两块相同的\emph{非理想}偏振片顺次放在自然光路中。
两片透振轴平行时的透射光强为正交时的
$\eta=10.0$ 倍。求 (a) 单块偏振片输出光的偏振度；
(b) 两片平行排列时输出光的偏振度。
\exsolution
设每片对沿主轴和正交轴的入射线偏振分量，
分别透射强度比例 $a,b$（$a\ge b\ge0$）。
对各轴原始强度各为 $I_0/2$ 的自然光，
\[
 I_\parallel=\frac{I_0}{2}(a^2+b^2),\qquad
 I_\perp=I_0 ab,
 \quad\eta=\frac{a^2+b^2}{2ab}.
\]
单片输出偏振度
\[
 P_1=\frac{a-b}{a+b},\qquad
 \boxed{P_1=\sqrt{\frac{\eta-1}{\eta+1}}
 =\sqrt{\frac9{11}}\simeq0.905}.
\]
双片平行排列后两主轴的强度比例分别为 $a^2,b^2$，
\[
 P_2=\frac{a^2-b^2}{a^2+b^2},\qquad
 \boxed{P_2=\sqrt{1-\frac1{\eta^2}}
 =\sqrt{0.99}\simeq0.995}.
\]
\excheck 单片结果约 $0.905$，双片约 $0.995$；
书末纯文本易漏根号，须以 $a,b$ 的代数定义判定。

\polarizationproblem{165}{214}
两束平行、等强度的线偏振光，其透振方向夹角很小，
为 $\phi$（见原图5.30）。
调整检偏棱镜透振方向，使出射两束光强相等，
有内角平分线 $A$ 与外角平分线 $B$ 两种位置。
棱镜均自相应位置微转 $\delta\phi\ll\phi$。
要求在 $A$ 位置引起的相对强度差变化
为 $B$ 位置的 $\eta=100$ 倍，求 $\phi$。
\exsolution
对强度为 $I$ 的入射光，
出射强度为 $I\cos^2\theta$。
在 $A$ 位置，原有两光强为
$I\sin^2(\phi/2)$（按原图 $A$ 为消光平分线），
两束的差对微转量的一阶项大小为
$2I\sin\phi\,|\delta\phi|$；
除以各束基准强度即
\[
 \left(\frac{\Delta I}{I_{\rm out}}\right)_A
 =4\cot\frac\phi2\,|\delta\phi|.
\]
在 $B$ 位置，以基准 $I\cos^2(\phi/2)$ 归一化，得到
\[
 \left(\frac{\Delta I}{I_{\rm out}}\right)_B
 =4\tan\frac\phi2\,|\delta\phi|.
\]
故
\[
 \eta=\cot^2\frac\phi2,\qquad
 \boxed{\phi=2\arctan\frac1{\sqrt\eta}\simeq11.4^\circ}.
\]
\excheck 相对\emph{两束强度之差}的一阶变化使系数为4；
只计算单束变化会少一倍，但比例 $\eta$ 不受影响。

\polarizationproblem{166}{214}
用菲涅耳公式证明：当光由折射率为 $n_1$ 的介质
入射折射率为 $n_2$ 的介质，且
$\tan\theta_1=n_2/n_1$ 时，
反射光完全偏振。此时反射光与折射光方向的夹角是多少？
\exsolution
平行入射面的电场分量的反射率为
\[
 R_p=\frac{\tan^2(\theta_1-\theta_2)}
 {\tan^2(\theta_1+\theta_2)}.
\]
布儒斯特条件
$\tan\theta_B=n_2/n_1$ 配合
$n_1\sin\theta_B=n_2\sin\theta_2$
给出 $\theta_B+\theta_2=90^\circ$，
故 $R_p=0$，反射光只含垂直入射面的电场分量。
反射线与折射线的夹角为
$180^\circ-(\theta_B+\theta_2)$，
于是
\[
 \boxed{\gamma=90^\circ}.
\]
\excheck 这是无损、非磁性各向同性介质的布儒斯特情形。

\polarizationproblem{167}{215}
自然光以布儒斯特角照射空气--玻璃界面。
用菲涅耳公式求 (a) 总反射率；
(b) 透射（折射）光的偏振度。
\exsolution
设玻璃折射率为 $n$（取表列近似 $n=1.50$）。
$p$ 分量在布儒斯特角不反射；
$s$ 分量反射率为
$R_s=((n^2-1)/(n^2+1))^2$。
自然光两线偏振分量各占一半，因此
\[
 \boxed{\rho=\frac12 R_s
 =\frac12\left(\frac{n^2-1}{n^2+1}\right)^2
 \simeq0.074}.
\]
透射两分量相对强度分别是
$1$ 和 $1-R_s$，于是
\[
 \boxed{P_t=\frac{R_s}{2-R_s}
 =\frac{\rho}{1-\rho}\simeq0.080}.
\]
\excheck 透射光在\emph{入射面内}的电矢量分量较强。

\polarizationproblem{168}{215}
自然光的强度为 $I_0$，以布儒斯特角入射水面。
已知反射\emph{能流}占 $\rho=0.039$，
求折射光束的强度（定义为垂直于折射光方向单位面积的能流率）。
\exsolution
能量守恒给出通过界面的能流率为
$(1-\rho)$ 倍入射能流率。
若入射光束横截面为 $A_i$，
界面光斑面积为 $A_i/\cos\theta_i$，
折射光束横截面为
$A_t=A_i\cos\theta_t/\cos\theta_i$。
因此
\[
 I_t=I_0(1-\rho)\frac{\cos\theta_i}{\cos\theta_t}.
\]
布儒斯特角满足
$\theta_i+\theta_t=\pi/2$、
$\tan\theta_i=n$，故
$\cos\theta_i/\cos\theta_t=1/n$。
取 $n_{\rm water}\simeq1.33$，
\[
 \boxed{I_t=\frac{1-\rho}{n}\,I_0
 \simeq0.722\,I_0}.
\]
\excheck 不得把透射总能流率比例 $1-\rho$
直接当作不同光束横截面的强度比例。

\polarizationproblem{169}{215}
线偏振光以布儒斯特角入射水面。
入射电矢量的振动方向与入射面成
$\phi=45^\circ$。求强度反射率。
\exsolution
将入射电矢量分为 $p,s$ 两分量。
在布儒斯特角 $R_p=0$，
仅 $s$ 分量反射：
\[
 \rho=R_s\sin^2\phi
 =\left(\frac{n^2-1}{n^2+1}\right)^2\sin^2\phi.
\]
以 $n=1.33,\ \phi=45^\circ$ 得
\[
 \boxed{\rho\simeq0.038}.
\]
\excheck 若 $\phi=0$（电矢量在入射面内），
反射率必须为零。

\polarizationproblem{170}{215}
窄束自然光以布儒斯特角射向厚透明平行平板，
上表面反射的能流比例 $\rho=0.080$。
求原图5.31中第1至第4束光各自的偏振度，
忽略相互重叠及光的吸收。
\exsolution
以 $p,s$ 分量分别追踪每次界面反射、透射。
布儒斯特角仅 $s$ 分量反射，
且自然光的总反射率为 $\rho$，
所以 $R_s=2\rho$。
第1、第3束都是界面反射束，故
\[
 \boxed{P_1=P_3=1}.
\]
第2束为一次透射束，偏振度
\[
 \boxed{P_2=\frac{1-(1-2\rho)}{1+(1-2\rho)}
 =\frac{\rho}{1-\rho}\simeq0.0870}.
\]
原图标出的第4束为经两界面透射的直行束，
两次透射使 $s$ 分量相对 $p$ 分量乘
$(1-2\rho)^2$，所以
\[
 \boxed{P_4=\frac{1-(1-2\rho)^2}{1+(1-2\rho)^2}
 =\frac{2\rho(1-\rho)}{1-2\rho(1-\rho)}
 \simeq0.173}.
\]
\excheck 第4束的 $P_4$ 应高于第2束的 $P_2$；
以强度透射系数追踪，不能对已是强度的系数再平方根。

\polarizationproblem{171}{215}
窄束单色光以布儒斯特角照射空气--玻璃平板
（原图5.31）；求 (a) 当入射光电矢量垂直于入射面时
直通透射束的强度；
(b) 当入射光为自然光时，直通透射束的偏振度。
忽略二次反射光束的再汇合。
\exsolution
在布儒斯特角，$p$ 分量界面反射率为0，
$s$ 分量为
\[
 R_s=\left(\frac{n^2-1}{n^2+1}\right)^2.
\]
透射束通过两界面，所以
\[
 \boxed{I_{\mathrm{out},s}
 =I_0(1-R_s)^2
 =I_0\frac{16n^4}{(1+n^2)^4}}.
\]
当玻璃 $n=1.50$ 时约为 $0.725\,I_0$。
自然光的两正交分量原强度各为 $I_0/2$；
$p$ 分量直通系数为1，
$s$ 分量直通系数为 $(1-R_s)^2$，故
\[
 \boxed{P=\frac{1-(1-R_s)^2}
 {1+(1-R_s)^2}\simeq0.159}.
\]
\excheck (a) 中印刷答案的分母幂有时在纯文本抽取时丢失，
以两次界面透射的 $(1-R_s)^2$ 为准。

\polarizationproblem{172}{215}
自然光以布儒斯特角穿过 $N$ 块厚的平行玻璃平板。
求 (a) 透射束的偏振度 $P_N$；
(b) 玻璃折射率取 $n=1.50$ 时，
$N=1,2,5,10$ 各对应何值。
\exsolution
定义 $\alpha=2n/(1+n^2)$。
单界面对 $s$ 分量的透射率
$1-R_s=\alpha^2$，
通过 $N$ 块平板共 $2N$ 个界面，
$s$ 分量透过比例 $\alpha^{4N}$；
$p$ 分量在布儒斯特角不受反射损失。
所以
\[
 \boxed{P_N=\frac{1-\alpha^{4N}}{1+\alpha^{4N}},
 \quad \alpha=\frac{2n}{1+n^2}}.
\]
代入 $n=1.50$，
\[
 \boxed{P_1\simeq0.16,\quad P_2\simeq0.31,\quad
 P_5\simeq0.67,\quad P_{10}\simeq0.92}.
\]
\excheck $N\to\infty$ 时 $P_N\to1$；
此处忽略平板内二次反射叠加和吸收。

\polarizationproblem{173}{215}
由菲涅耳方程求 (a) 自然光垂直入射空气--玻璃界面的反射率；
(b) 一束近轴自然光穿过同轴五片玻璃透镜时，
仅因各界面反射损失的能流比例（忽略二次反射）。
\exsolution
法向入射两偏振的反射率相同，
\[
 \boxed{\rho=\left(\frac{n-1}{n+1}\right)^2}.
\]
取 $n=1.50$，$\rho=0.040$。
五片透镜有十个空气--玻璃界面，
每次仅保留前向透射的 $(1-\rho)$ 部分，
于是
\[
 \boxed{\frac{\Delta\Phi}{\Phi_0}
 =1-(1-\rho)^{10}
 =1-0.96^{10}\simeq0.335}.
\]
\excheck 这是纯\emph{反射}损失，未计吸收和表面散射。

\polarizationproblem{174}{215}
玻璃表面覆盖一层折射率为 $n'$ 的透明薄膜，
玻璃折射率为 $n$，外部为空气。
证明：忽略二次反射时，
若来自薄膜两个界面的反射光振幅大小相等，
需满足何种折射率关系？
\exsolution
法向入射时，空气--膜界面的振幅反射系数为
\[
 r_{01}=\frac{1-n'}{1+n'},
\]
膜--玻璃界面为 $r_{12}=(n'-n)/(n'+n)$。
在 $1<n'<n$ 时两者同号，等振幅条件写为
$|r_{01}|=|r_{12}|$；
解方程
\[
 \frac{n'-1}{n'+1}=\frac{n-n'}{n+n'}
 \quad\Longrightarrow\quad (n')^2=n.
\]
故
\[
 \boxed{n'=\sqrt n}.
\]
\excheck 当法向入射、薄膜无吸收且忽略高阶反射时这是经典
增透膜等振幅条件；原书 PDF 的公式文字层可能漏掉根号，
原印字形应再复核。若严格计入膜界面透射振幅系数，
则应采用完整薄膜振幅公式，不把此一阶近似充当精确结论。

\polarizationproblem{175}{216}
自然光以 $\theta_i=45^\circ$ 入射玻璃表面。
用菲涅耳方程求 (a) 反射光偏振度；
(b) 折射光偏振度。取玻璃折射率 $n=1.50$。
\exsolution
斯涅耳定律给
$\theta_t=\arcsin(\sin45^\circ/n)\simeq28.13^\circ$。
两种线偏振的强度反射率为
\[
 R_s=\frac{\sin^2(\theta_i-\theta_t)}
 {\sin^2(\theta_i+\theta_t)},\qquad
 R_p=\frac{\tan^2(\theta_i-\theta_t)}
 {\tan^2(\theta_i+\theta_t)}.
\]
数值约为 $R_s=0.0920,\ R_p=0.00847$。
反射光偏振度及透射光偏振度的模分别为
\[
 \boxed{P_r=\frac{R_s-R_p}{R_s+R_p}\simeq0.83},
 \qquad
 \boxed{P_t=\frac{R_s-R_p}{2-R_s-R_p}\simeq0.044}.
\]
\excheck 两光束主要偏振方向不同：
反射光主要为 $s$ 分量，透射光主要为 $p$ 分量。
